\documentclass[11pt,letterpaper]{article}
% Shared layout. Compile each document from this directory.
\usepackage[margin=0.8in,headheight=14pt,headsep=16pt,footskip=22pt]{geometry}
\usepackage[T1]{fontenc}
\usepackage{lmodern,amsmath,amssymb,booktabs,enumitem,titlesec,fancyhdr,microtype,xcolor,hyperref}
\definecolor{inkblue}{RGB}{25,65,100}
\hypersetup{colorlinks=true,linkcolor=inkblue,urlcolor=inkblue,filecolor=inkblue,pdfauthor={AI-assisted course study notes}}
\setlength{\parindent}{0pt}
\setlength{\parskip}{5pt}
\setlength{\emergencystretch}{1.5em}
\setlength{\abovedisplayskip}{7pt}
\setlength{\belowdisplayskip}{7pt}
\titleformat{\section}{\normalsize\bfseries\color{inkblue}}{}{0pt}{}
\titlespacing*{\section}{0pt}{9pt}{3pt}
\setlist[itemize]{leftmargin=15pt,itemsep=2pt,topsep=3pt,parsep=0pt}
\setlist[enumerate]{leftmargin=17pt,itemsep=3pt,topsep=3pt,parsep=0pt}
\pagestyle{fancy}
\fancyhf{}
\fancyhead[L]{\small ECE 750T\enspace /\enspace Responsible AI}
\fancyhead[R]{\small Fall 2026\enspace /\enspace Study notes}
\fancyfoot[R]{\small\thepage}
\renewcommand{\headrulewidth}{0.3pt}
\newcommand{\lesson}[2]{%
  \hypersetup{pdftitle={ECE 750T: #1 -- #2}}%
  {\LARGE\bfseries #1}\par\vspace{3pt}
  {\large #2}\par\vspace{3pt}
  {\footnotesize AI-assisted reconstruction from course slides, not a lecture transcript.
  Discussion reasoning is inferred; classroom responses were not supplied.}\par
  \vspace{3pt}\hrule\vspace{5pt}}
\newcommand{\sources}[1]{\par\vspace{3pt}{\footnotesize\textbf{Slide trail.} #1
  Page numbers refer to PDF pages, not the slides' continuous printed numbering.}\par}
\newcommand{\Prob}{\mathbb{P}}
\newcommand{\E}{\mathbb{E}}
\newcommand{\Yhat}{\widehat{Y}}
\newcommand{\indep}{\mathrel{\perp\!\!\!\perp}}

\begin{document}
\lesson{Lecture 4}{Measurement and data curation: precision is not validity}

\textbf{The lesson.} More precise estimates of a recorded quantity do not establish
that the quantity represents the thing we care about. The lecture connects
sampling error and model complexity to construct validity, biased observation,
feedback under optimization, and documentation that supports actual use decisions.

\section*{What more data improves, and under which assumptions}
For a fixed classifier evaluated on $n$ independent draws from the target
population, correctness is Bernoulli with parameter $p$:
\[
\E[\widehat p]=p,\qquad
\mathrm{SE}(\widehat p)=\sqrt{p(1-p)/n}.
\]
Halving the standard error requires four times as many examples. At $p=.8$,
$n=25,100,400$ give SEs $.08,.04,.02$. This is an average precision statement:
an individual added sample can move the estimate farther from the true value.

For a fixed, finite hypothesis class $H$, i.i.d. examples, and 0--1 loss, with
probability at least $1-\delta$,
\[
\sup_{h\in H}|R(h)-\widehat R(h)|
\le\sqrt{\frac{\log(2|H|/\delta)}{2n}}.
\]
More candidate models enlarge the bound; more data shrink it. The simultaneous
guarantee differs from estimating one model fixed in advance. Regularization
controls complexity, but this particular finite-class bound cannot be applied
literally to an unrestricted neural-network parameter space. Optimizing a smooth
surrogate such as cross-entropy is another distinction: the training objective
and the reported accuracy are different quantities.

\section*{Observation and experimental design change the inference}
The historical examples show data collection as purposeful work for understanding
populations and acting on social conditions.

The returning-aircraft illustration is \textbf{selection/survivorship bias}:
damage is observed conditional on returning. It misses aircraft that did not
survive. In the fertilizer example, soil quality causes both fertilizer use and
yield, so the observed association mixes the treatment effect with a
\textbf{confounder}. Increasing the same biased observational sample does not
by itself resolve either problem.

Blocking agricultural plots by soil conditions puts every treatment within
comparable conditions; randomizing treatment positions within blocks reduces
systematic assignment differences. This explains why design matters before
estimation and why a correlation does not identify the desired causal effect.

\section*{Separate construct, measurement, and objective}
In the health example, \emph{need} is the construct, recorded spending is the
measurement used as a target, and prediction loss is the optimized objective.
Spending also reflects prices, access, and service use. Better cost prediction
can therefore sharpen a poor proxy for unmet need.

Construct validity concerns whether evidence and theory support interpreting a
measurement as the intended construct for a particular use. Reliability concerns
consistency. A timed coding test may consistently rank speed while omitting
debugging and collaboration; accurate spending records can omit unmet care.
More data or repeated measurements cannot alone establish validity.

\newpage
\section*{Biased labels can change the fairness conclusion}
The COMPAS discussion separates recorded rearrest $Y$ from reoffending $Y^*$.
Detection and enforcement sit between them. If that observation process differs
across groups, fairness measured against $Y$ need not match fairness against
$Y^*$. Varying assumptions about unrecorded offenses in a sensitivity analysis
tests whether the conclusion changes; it does not reveal the unobserved truth.

Predictive policing makes a feedback mechanism visible: recorded arrests reflect
police activity as well as underlying drug use. Sending officers where arrests
were recorded can produce further records there. The target can partly measure
the institution's past actions. The healthcare case similarly shows why equal
treatment conditional on a recorded score can coexist with unequal unmet need.

\section*{Optimization can break a proxy's usefulness}
Goodhart's law warns that a statistical relationship can change when a measure
becomes a policy target. People or agents adapt to the incentive. In the boat-race
example, points are awarded for hitting targets; the agent repeatedly circles
through respawning targets instead of finishing the race. Reward rises while the
intended outcome fails: \emph{reward hacking}.

Identify the exploitable gap between measure and goal, then test a redesign under
optimization. More examples of high reward do not establish success. The same
reasoning applies to benchmarks repeatedly used for model selection.

\section*{Documentation is evidence for a use decision}
Data pass through collection, annotation, training, and deployment. Assumptions
about population, label meaning, missing values, and collection conditions can
be lost at these handoffs, producing downstream data cascades. The component
datasheet analogy asks engineers to specify operating conditions and limitations,
as they would for a sensor.

\textbf{Datasheets} describe dataset purpose, composition, collection, uses, and
maintenance. \textbf{Model cards} describe model uses, evaluation, and limitations.
\textbf{System cards} address a system's capabilities, risks, and mitigations.
They should help a reader assess suitability for a specified use.

The smile-detector example proposes selecting photographs of children. Evaluation
on CelebA by annotated age and gender does not establish performance for children.
A useful revision requests relevant evaluation images, label definitions, and errors
by age and capture conditions. Missing evidence establishes uncertainty, not an
invented failure result. This use--evidence--gap--revision chain is central to A2.

\section*{Reproduction answers a further, distinct question}
Code, provenance, evaluation procedures, and reporting checklists support
reproduction. Different seeds can change results with fixed algorithms and
hyperparameters. For $m$ independent runs, $s/\sqrt m$ estimates the mean's SE;
it differs from individual-run SD and test-example uncertainty. Reproducing a score
does not establish validity or deployment usefulness.

\textbf{Midterm emphasis (inference).} Know the fourfold-sample calculation and
the interpretation and assumptions of the complexity bound. Strong conceptual
priorities are validity versus reliability, rearrest versus reoffending, confounding
versus selection, reward hacking, and a concrete documentation improvement with
the evidence it needs. These remain in scope despite not appearing on the reference sheet.

\sources{\href{../L04.pdf}{L04.pdf}, pp.~2--6 (estimation/complexity),
10--20 (design, validity, labels, Goodhart), 21--29 (documentation/reproduction);
annotated copy consulted. A2: pp.~1--2.}
\end{document}
